% =====================================================================
%  Thesis assignment (consolidated) — ENGLISH (source of truth)
%  Consolidated from the project meetings docs/termine/ (Termin 1-9) AND the
%  full architecture docs (docs/architektur/14 organ metaphor, bausteine
%  matrix N, kapitel 01/04). German counterpart: de.tex.
% =====================================================================
{\large\textbf{Thesis Assignment}}\\[0.4em]
{\itshape Active Cache-Aware Hardware Adaptation:\\ A Cache Engine for Trie-Based Index Structures}

\bigskip
\minisec{Context and motivation}
Trie-based index structures are central to modern in-memory database systems,
and the Chair of Databases focuses explicitly on hardware-conscious optimisation
of such systems. Since the Adaptive Radix Tree, numerous cache-aware
techniques---node layout, prefix compression, prefetching, succinct
representations, and allocation---have been proposed, yet each publication uses
its own data structures, workloads, and measurement methodology. No common
framework decomposes these designs into interchangeable building blocks and
measures them on a single, fair basis.

\minisec{Objective}
The thesis shall design and prototype an axis-library framework (a
\emph{cache engine}) that decomposes cache-aware trie indexing into orthogonal,
independently exchangeable building-block axes, reconstructs state-of-the-art
algorithms as explicit axis configurations, and measures both individual axes
and whole algorithms on a common, CPU-only basis. As a probe for
state-of-the-art membership, the thesis contributes \mbox{PRT-ART} (a Probst
Redirect Tree~/ ART hybrid).

\minisec{Methodical approach: algorithms as ``animals'', axes as ``organs''}
The framework rests on an anatomical metaphor. A search algorithm is treated as
an \emph{animal} (organism); the sub-tasks that \emph{every} algorithm must
perform are its \emph{organs}. Each organ is one orthogonal \emph{axis}: just as
every ruminant has a stomach but cow, deer, and sheep realise it differently,
every search algorithm has a node-layout organ, a traversal organ, an allocation
organ, and so on, each in a different expression. A concrete algorithm is then
\emph{one composition} --- a single point in the Cartesian product of all axes;
a state-of-the-art design (ART, HOT, Masstree, \ldots) is exactly one such
point. Permuting the axes is a \emph{genetic experiment}: exchanging organs
between animals yields new composite algorithms with previously unknown
properties. The research goal is to find the common \emph{anatomy} of a search
algorithm and, within the resulting permutation space, the fastest ``animal''
--- the optimal cache-line-aware composition --- for a given data-type-favourable
load pattern.

\minisec{Building-block axes and cache-line relevance}
A task of the thesis is to determine which axes (organs) every search algorithm
exposes. The consolidated matrix comprises \textbf{fourteen main axes plus about
thirty sub-axes}: (1)~Page-Type, (2)~Node-Type, (3)~Traversal --- split into
3.A~algorithm-side, 3.B~cache-memory-side and 3.M~mapping ---, (4)~ValueHandle,
(5)~Memory-Layout, (6)~Allocator (allocation, reclamation, NUMA, huge-page,
free-list), (7)~Prefetch, (8)~Concurrency (pattern, locking-mode), (9)~ISA,
(10)~Measurement, (11)~Telemetry, (12)~Hardware-Strategy, (13)~Scheduling-Strategy,
and (14)~Identity. The \textbf{cache-line-aware-relevant} organs are in
particular: Node-Type and Memory-Layout (how nodes occupy 64-byte cache lines),
the \emph{cache-memory traversal} (3.B) together with its \emph{mapping} (3.M,
which translates the developer's non-cache-coherent algorithm view into the
actual cache mechanics on the memory side), the Allocator (NUMA affinity and
huge pages govern cache/TLB locality), Prefetch, ISA, and Hardware-Strategy.

\minisec{Why exhaustive permutation is required}
Cache-line-aware behaviour does not arise from any single organ but from the
\emph{interaction} of many --- node layout $\times$ memory layout $\times$
allocator $\times$ prefetch $\times$ traversal mapping $\times$ ISA. An organ
cannot be tuned for cache lines in isolation; the cross-effects between axes are
decisive. The thesis therefore measures the configuration space \emph{fully and
permutatively}, with no reduction (``measurement always as complete as
possible''); the Cartesian product spans more than~$10^{11}$ permutations, which
the builder emits as individual, ABI-stable measurement binaries.

\minisec{Sub-tasks}
\begin{enumerate}
  \item Systematise the organs of cache-aware trie indexes into a canonical
        permutation matrix of orthogonal axes and identify the cache-line-aware-relevant
        ones.
  \item Reconstruct state-of-the-art designs (ART, HOT, Masstree, CoCo-trie,
        START, B\textsuperscript{2}-tree, Wormhole, SuRF, and further Tier-2/3
        designs) as explicit axis compositions (reference ``animals'').
  \item Design and implement PRT-ART as a probe with its own organs (redirect
        nodes, adaptive local search pages, external value arenas,
        cache-line-aware layout).
  \item Build a reproducible measurement pipeline
        (binary~$\to$~CSV~$\to$~LaTeX~$\to$~diagram) with profile-aware YCSB
        workload routing and hardware-counter collection.
  \item Conduct the three mandatory measurement series exhaustively: (A)~probe
        vs.\ state of the art, (B)~cache-engine permutations
        (allocator~$\times$~layout~$\times$~prefetch~$\times$~\ldots),
        (C)~merge/regression of old vs.\ new variants.
  \item Identify, from the full permutation space, the fastest cache-line-aware
        compositions per data-type-favourable load pattern (research goal).
\end{enumerate}

\minisec{Method and evaluation}
The evaluation is CPU-only and compares the compositions against established
baselines under varied data volumes, key lengths, prefix overlap, local density,
skew, hit/miss ratios, and value sizes (the data-type-favourable load patterns).
Measured quantities include exact/prefix lookup, prefix enumeration,
insert/update/delete, p50/p95/p99 latency, throughput, bytes per key, and
hardware counters (cycles, instructions, L1/L2/LLC and dTLB misses, branch
mispredictions). Single-thread behaviour is studied first, followed by
read-parallel multi-thread scenarios.

\minisec{Scope}
The core contribution is CPU-only and single-thread with read scaling. GPU is
out of scope; write parallelism, SIMD, and engine/optimiser integration are
optional extensions.
